\paragraph{有限状态实现下的时间像与离散化}

\begin{proposition}[有限状态实现下的时间像：离散量子化与稠密化二分]\label{prop:xi-finite-state-time-image-discrete-or-dense}
设可见群胚 \(\mathcal G\) 具有有限生成的有向实现（有限状态、有限标记边集 \(E\)），并给定可加时间余圈 \(\widetilde{\ell}:\mathcal G\to\mathbb R\) 满足 \(\widetilde{\ell}(gh)=\widetilde{\ell}(g)+\widetilde{\ell}(h)\)。
若生成边上的取值集合 \(\{\widetilde{\ell}(e):e\in E\}\) 有限，则
\[
\widetilde{\ell}(\mathcal G)\subseteq
H:=\Bigl\langle \widetilde{\ell}(e):e\in E\Bigr\rangle_{\mathbb Z}\ \le\ (\mathbb R,+).
\]
并且 \(H\) 作为 \(\mathbb Z\)-模有限生成，故满足二分律：
\[
H=h\mathbb Z\ \text{（离散）}\qquad\text{或}\qquad \overline{H}=\mathbb R\ \text{（稠密）}.
\]
在离散情形中，\(h\) 可取为 \(H\cap(0,\infty)\) 的最小正元（等价于 \(\inf(H\cap(0,\infty))\)），并有 \(\widetilde{\ell}(\mathcal G)\subset h\mathbb Z\)。
此外，\(\widetilde{\ell}(\mathcal G)\) 至多可数，因而任何真正连续的时间像（不可数）只能通过状态数趋于无穷的极限获得。
\end{proposition}

\begin{proof}
因 \(\mathcal G\) 由有限边集生成，任意 \(g\in\mathcal G\) 的 \(\widetilde{\ell}(g)\) 为生成边取值的整数线性组合，故 \(\widetilde{\ell}(\mathcal G)\subseteq H\)。
而任意 \(\mathbb Z\)-有限生成子群 \(H\le\mathbb R\) 要么含有最小正元（从而必为该最小正元生成的循环群 \(h\mathbb Z\)），要么在 \(\mathbb R\) 中稠密；这是加法子群的标准分类。
最后，有限生成群 \(H\cong \mathbb Z^k\) 可数，故 \(\widetilde{\ell}(\mathcal G)\) 可数。
\end{proof}

\paragraph{欧拉乘积、Holonomy 与 RH 的同调化归约}

\begin{theorem}[时间欧拉分解定理：混合曲率零强制配分函数的素元乘积分解]\label{thm:xi-time-euler-product-under-zero-mixed-curvature}
设 \((\mathcal M_\times,\circ_\times)\subset \mathrm{Mor}(\mathbf{PW})\) 为乘法复合幺半群，并假设存在时间余圈 \(\widetilde{\ell}\) 对 \(\circ_\times\) 严格可加。
令 \(\mathcal P\subset\mathcal M_\times\) 为 \(\circ_\times\)-不可约元（素元）集合。
若满足：
\begin{enumerate}
  \item \(\mathcal M_\times\) 为交换、左消去且为素元的自由交换幺半群（等价地，每个元具有唯一的素元分解，按重数计）；
  \item 加乘交错方块的混合曲率恒为零（定理 \ref{thm:xi-mixed-curvature-zero-criterion}），从而时间在加/乘两层间无中心补偿缺陷；
\end{enumerate}
则对任意处于绝对收敛域的 \(\lambda>0\)，乘法配分函数
\[
Z_\times(\lambda):=\sum_{w\in\mathcal M_\times}e^{-\lambda \widetilde{\ell}(w)}
\]
具有欧拉乘积分解
\[
Z_\times(\lambda)=\prod_{p\in\mathcal P}\bigl(1-e^{-\lambda \widetilde{\ell}(p)}\bigr)^{-1}.
\]
反之，若对某开区间 \(\lambda\in(\lambda_0,\infty)\) 上 \(Z_\times\) 具有上述欧拉乘积且 \(\mathcal P\) 生成 \(\mathcal M_\times\)，则混合曲率必为零（等价于加法更新与乘法更新在时间意义下可交换且无需额外中心补偿）。
\end{theorem}

\begin{proof}
在唯一素元分解与可加性下，
\(
\widetilde{\ell}\bigl(\prod_p p^{n_p}\bigr)=\sum_p n_p\widetilde{\ell}(p)
\)。
对 \(\lambda\) 处于绝对收敛域时，按素元分解改写求和并交换求和与乘积，得到
\[
Z_\times(\lambda)
=\prod_{p\in\mathcal P}\sum_{n\ge 0}e^{-\lambda n\widetilde{\ell}(p)}
=\prod_{p\in\mathcal P}(1-e^{-\lambda\widetilde{\ell}(p)})^{-1}.
\]
反向方向中，欧拉乘积在一段开区间上成立迫使各阶系数满足素元独立的乘法型生成结构，从而与“交错方块无缺陷”的路径无关性相容；若混合曲率非零，则加/乘交错会引入不可由边界吸收的中心补偿项，破坏该乘法型分解，故必须有 \(K\equiv 0\)。
\end{proof}

\begin{theorem}[时间统一推出 RH 的同调化归约：偏离临界线等价于非平凡时间 Holonomy]\label{thm:xi-rh-cohomological-reduction-holonomy}
假设存在完成化表示，使得经典完成化函数 \(\Xi(s)\) 可写为时间加权行列式族的限制：
\[
\Xi(s)=\det\!\bigl(I-e^{-s}\mathcal L_{\theta(s)}\bigr),
\qquad
\theta(s):=\Re(s)-\tfrac12,
\]
其中 \(\theta\mapsto\mathcal L_{\theta}\) 为由可审计局部演化与时间余圈诱导的转移算子族，并满足时间反演共轭
\(
\mathcal L_{\theta}\sim \mathcal L_{-\theta}
\)
对应于 \(\Xi(s)=\Xi(1-s)\)。
若进一步满足：
\begin{enumerate}
  \item 严格正时间：任意非平凡闭环的时间和非零；
  \item 尺度--时间联络的 Holonomy 群 \(\mathfrak H\) 平凡（定理 \ref{thm:xi-scale-time-connection-holonomy-criterion}）；
  \item 临界切片的压力判据成立（定理 \ref{thm:xi-critical-slice-pressure-zero-drift-lyapunov}），从而 \(\theta=0\) 为唯一的“双零”切片；
\end{enumerate}
则对任意 \(\theta\neq 0\) 必有 \(\rho(\mathcal L_{\theta})\neq 1\)，从而 \(\det(I-e^{-s}\mathcal L_{\theta})\neq 0\) 不可能发生在 \(\theta\neq 0\) 的切片上。
因此所有非平凡零点必满足 \(\theta(s)=0\)，即 \(\Re(s)=\tfrac12\)。
在上述表示假设下，RH 可被等价重述为“尺度--时间 Holonomy 消失”所强制的临界切片唯一性。
\end{theorem}

\begin{proof}
由定理 \ref{thm:xi-critical-slice-pressure-zero-drift-lyapunov}，在可审计微可逆与时间反演口径下，\(\theta=0\) 为唯一同时满足零漂移与零 Lyapunov 的切片。
Holonomy 平凡排除在跨尺度拼合中引入额外中心补偿的自由度，从而临界性不再可沿 \(\theta\neq 0\) 方向保持。
因此对 \(\theta\neq 0\) 必有 \(\rho(\mathcal L_{\theta})\neq 1\)，与行列式零点发生在非临界谱条件矛盾；结论遂得。
\end{proof}

\begin{proposition}[RH 的有限可判据：仅需检查有限生成的时间曲率矩阵的整秩消失]\label{prop:xi-rh-finite-curvature-matrix-criterion}
在定理 \ref{thm:xi-rh-cohomological-reduction-holonomy} 的表示假设下，若协议来自有限生成的收敛重写系统，则 Holonomy 群 \(\mathfrak H\) 由有限多个“尺度--时间临界方块”的缺陷生成。
将这些缺陷在某一中心坐标系（\(\mathbb Z^d\)）下记为整数向量，并排列成矩阵 \(B\in\mathbb Z^{r\times d}\)（行对应临界方块，列对应中心基向量），则
\[
\mathfrak H\cong \mathrm{im}(B)\subset\mathbb Z^d,
\qquad
\mathrm{rank}(\mathfrak H)=\mathrm{rank}_{\mathbb Z}(B).
\]
特别地在标量缺陷 \(d=1\) 情形，\(\mathfrak H=\{0\}\) 当且仅当 \(B=0\)。
因此在该框架内，“Holonomy 消失”归约为对有限整数矩阵的零性/秩判别，且与 gauge 选择无关。
\end{proposition}

\begin{proof}
收敛重写系统的所有高阶重叠由有限临界对生成；Holonomy 作为交换方块的缺陷累积，故由有限临界方块的缺陷生成。
把生成元在 \(\mathbb Z^d\) 基下坐标化即得矩阵描述；秩结论由 \(\mathbb Z\)-线性代数给出。
\end{proof}

\paragraph{正时间规范、Smith 归约与 PBW 可交换化}

\begin{proposition}[正时间的同调正锥判据：全边正 gauge 的存在性]\label{prop:xi-positive-time-cone-criterion}
设 \(\Gamma=(V,E)\) 为有限强连通有向图，\(t:E\to\mathbb R\) 为时间权，并作 gauge 变换 \(t\mapsto t+\partial b\)（\(b:V\to\mathbb R\)）。
则存在 \(b\) 使得 \(t+\partial b\) 在所有边上严格为正，当且仅当对每条有向闭环 \(\gamma\) 有
\[
\sum_{e\in\gamma} t(e)>0.
\]
等价地，\([t]\in H^1(\Gamma;\mathbb R)\) 落在由所有有向闭环定义的正锥
\[
\mathcal C_+:=\Bigl\{[c]\in H^1:\ \langle[c],[\gamma]\rangle>0\ \text{对一切非平凡闭环 }\gamma\Bigr\}
\]
的内部。
\end{proposition}

\begin{proof}
对任意闭环 \(\gamma\)，有 \(\sum_{e\in\gamma}\partial b(e)=0\)，故闭环和是同调不变量。
若存在 \(b\) 使各边正，则任意闭环和必正，得必要性。
反向为差分约束的可行性判据：不等式 \(t(e)+b(r(e))-b(s(e))>0\) 等价于严格差分约束系统，其可行性当且仅当不存在闭环总权 \(\le 0\)；在本设定中即为所有闭环和严格为正。
\end{proof}

\begin{proposition}[时间缺陷群的 Smith 归约：中心扩张的自由度与扭结结构可被完全计算]\label{prop:xi-time-defect-smith-reduction}
沿命题 \ref{prop:xi-rh-finite-curvature-matrix-criterion} 的矩阵化描述，设缺陷矩阵为 \(B\in\mathbb Z^{r\times d}\) 并令 \(\mathfrak H:=\mathrm{im}(B)\subset\mathbb Z^d\)。
则存在可逆整数矩阵 \(U\in\mathrm{GL}_r(\mathbb Z)\)、\(V\in\mathrm{GL}_d(\mathbb Z)\) 使
\[
UBV=\mathrm{diag}(s_1,\dots,s_k,0,\dots,0),
\qquad
s_i\mid s_{i+1}.
\]
从而缺陷余商具有标准分解
\[
\mathbb Z^d/\mathfrak H
\ \cong\
\bigoplus_{i=1}^k\mathbb Z/s_i\mathbb Z\ \oplus\ \mathbb Z^{d-k}.
\]
其中自由秩 \(d-k\) 给出最小中心扩张的连续参数维数，扭因子 \(s_i\) 给出不可消去的离散“时间扭结”。
\end{proposition}

\begin{proof}
由 Smith 标准形定理，任意整数矩阵在左右可逆整数变换下可化为上述对角形；对角元给出像子群的结构不变量，从而导出商群分解。
\end{proof}

\begin{theorem}[时间滤过的 PBW 型定理：混合曲率零当且仅当关联分级代数可交换]\label{thm:xi-time-filtration-pbw-graded-commutative}
令 \(\mathbb k[\mathbf{PW}]\) 为协议范畴代数（\(\mathbb k\) 为特征零域），并以时间长度 \(\widetilde{\ell}\) 定义滤过
\[
F_t:=\mathrm{span}\{w:\widetilde{\ell}(w)\le t\}.
\]
若 \(\widetilde{\ell}\) 对允许复合严格可加，则 \(\mathrm{gr}_{\widetilde{\ell}}\mathbb k[\mathbf{PW}]\) 为分级代数。
设系统存在两类生成（对应加法层与乘法层）并形成交错方块，则
\[
\mathrm{gr}_{\widetilde{\ell}}\mathbb k[\mathbf{PW}]\ \text{为交换代数}
\quad\Longleftrightarrow\quad
K\equiv 0,
\]
其中 \(K\) 为定理 \ref{thm:xi-mixed-curvature-zero-criterion} 的混合曲率。
\end{theorem}

\begin{proof}
若 \(K\equiv 0\)，交错方块两路径在时间主项上无缺陷，故交换关系在关联分级的首项中成立，从而分级代数可交换。
反之，若关联分级可交换，则所有交错生成元在首项中交换，故交换方块的时间缺陷必须为零，即 \(K\equiv 0\)。
\end{proof}

\paragraph{时间边界与统计退化}

\begin{theorem}[时间边界的 Patterson--Sullivan 密度：KMS 平衡态与边界共形测度等价]\label{thm:xi-time-boundary-patterson-sullivan-kms}
若时间度量 Cayley 图满足线性等周界（定理 \ref{thm:xi-audit-dehn-function-time-geometry} 的双曲性条件），并设增长指数 \(\kappa\) 存在。
对 \(\lambda>\kappa\) 定义 Poincar\'e 级数
\[
\mathcal P_\lambda(x,y):=\sum_{g:x\to y}e^{-\lambda \widetilde{\ell}(g)}.
\]
则存在唯一（至尺度）\(\lambda\)-共形的边界密度族 \(\{\nu_x^{(\lambda)}\}\) 支持于 Gromov 边界 \(\partial\)，并满足由 Busemann--时间函数给出的 Radon--Nikodym 关系。
并且该边界密度诱导的极端态与由 Gibbs 权重 \(e^{-\lambda\widetilde{\ell}}\) 所定义的 KMS 平衡态之间存在自然一一对应。
\end{theorem}

\begin{proof}
在双曲度量下，Poincar\'e 级数收敛域与共形密度的存在性由 Patterson--Sullivan 构造给出，且增长指数为其临界参数。
KMS 态与共形密度的对应是由时间余圈诱导的动力系统（模流/时间流）在边界上的共形性所给出的标准对应；极端分解与边界测度的极端性一致。
\end{proof}

\begin{proposition}[缺陷秩控制大偏差退化：\(\mathrm{rank}(\mathfrak H)\) 等于速率函数的仿射平坦维数]\label{prop:xi-defect-rank-controls-ldp-flatness}
设压力函数 \(P(s)=\log\rho(A_s)\) 在 \(s>s_c\) 可微且严格凸，其中 \(A_s\) 为时间加权转移算子。
对 gauge 扰动 \(\widetilde{\ell}\mapsto \widetilde{\ell}+\partial b\) 有 \(P\) 不变；对更一般的中心缺陷方向（由 Holonomy 缺陷群 \(\mathfrak H\) 给出）压力沿某些线性方向保持不变。
令大偏差速率函数 \(I\) 为 \(P\) 的 Legendre 对偶，则
\[
\dim\bigl(\text{\(I\) 的最大仿射平坦子集}\bigr)=\mathrm{rank}(\mathfrak H).
\]
特别地，\(\mathfrak H=\{0\}\) 当且仅当 \(I\) 严格凸；一旦存在非平凡缺陷自由度，速率函数必出现相应维数的平坦“相共存”方向。
\end{proposition}

\begin{proof}
压力在缺陷方向上的不变性意味着 \(P\) 在相应线性子空间上具有平坦方向；Legendre 对偶将平坦方向对应为速率函数的仿射平坦面。
维数由该不变子空间（等价于缺陷自由度）的秩给出。
\end{proof}

\paragraph{周期和的有限判别性：同调类的完全重建}

\begin{proposition}[周期和决定时间同调类：有限基循环的时间和相等即余边界等价]\label{prop:xi-cycle-sums-determine-time-cohomology}
设 \(\Gamma=(V,E)\) 为强连通有限有向图，\(t,t':E\to h\mathbb Z\) 为两组时间权，视为 \(1\)-余圈。
则
\[
t-t'=\partial b\ \text{对某 }b:V\to h\mathbb Z
\quad\Longleftrightarrow\quad
\sum_{e\in\gamma}t(e)=\sum_{e\in\gamma}t'(e)\ \text{对所有有向闭环 }\gamma.
\]
并且“对所有闭环”的检验可缩减为对任意生成 \(H_1(\Gamma;\mathbb Z)\) 的有限基本循环族即可：取任一生成树 \(T\)，对每条弦边 \(e\in E\setminus T\) 所形成的基本循环 \(\gamma_e\) 构成整数循环空间一组基，检查有限多条 \(\{\gamma_e\}\) 的时间和相等即推出同调等价。
\end{proposition}

\begin{proof}
必要性由闭环上 \(\sum\partial b=0\) 直接得出。
反向方向是有限图上 \(1\)-上同调的标准刻画：若两余圈在一组循环基上的配对一致，则其差在 \(H^1\) 中为零，从而为余边界。
\end{proof}
